\begin{description}
\item[(13.1)] A single electron deposits energy in the CCD, as follows:
\[ E_{\mathrm{deposited}}
   = \text{stopping power} \times \rho_{\mathrm{si}} \times \text{thickness} \]
\[ E_{deposited}
   = 3.012\,\frac{MeV\,cm^2}{g} \times \frac{2.34\,g}{cm^3}
     \times 0.06\,cm = 422.9~keV \]
\hfill(4 points)

Since an electron with 15\,MeV energy deposits 422.9\,keV, we must calculate
the number of pairs electron/hole
\[ 422.9\,keV \times \frac{1}{2.36\,eV}
   = 1.79 \times 10^{5}\,\frac{e}{h} \]
\hfill(4 points)

How many pixels will be able to excite $1.79 \times 10^{5}$ pairs of e/h?
\[ 1.79 \times 10^{5} \times \frac{1}{250} = 716~pixels \]
\hfill(2 points)

\item[(13.2)] The number of electrons entering the detector area is:
\[ flux \,\times\, area \,\times\, t \]
\hfill(1 point)
\[ \frac{600\,e}{cm^2\,s} \times [1.3 \times 1.3]\,cm^2 \times 0.03\,s
   = 30\,e^{-} \]
\hfill(1 point)

$\sim$30 electrons enter the detector area, and we know that one electron
excites 716 pixels, the total number of excited pixels is:
\[ 716\,\frac{pixels}{electron} \times 30.42~electrons
   = 2.18 \times 10^{4}\,pixels \]
\hfill(1 point)

If the CCD has a total of $1024 \times 1024$ pixels
$= 1.048 \times 10^{6}\,pixels$ then the total percentage of pixels that
will be driven in a single image is $\sim 2.08\%$. \hfill(2 points)
% sic: "driven" (presumably "excited") as printed in the source.
\end{description}
